\begin{figure}[!b]
\centering
\begin{tikzpicture}[>=Stealth,thick,font=\small]
  \coordinate (n) at (0,0);
  \draw[linkblue] (-4.6,0.2) .. controls (-2.7,0.8) and (-1.3,0.5) .. (n);
  \draw[linkblue] (n) .. controls (1.3,-0.5) and (2.7,-0.8) .. (4.6,-0.2);
  \fill (n) circle (2pt);
  \node[above=5pt] at (-3.2,0.45) {$(C_1,x_1)$};
  \node[below=5pt] at (3.2,-0.45) {$(C_2,x_2)$};
  \node[above=7pt] at (n) {$B\mu_{b_r}$};
  \node[align=center] at (-2.8,-0.75)
    {$u\mapsto\zeta u$\\$u_{1,r}(\zeta)$};
  \node[align=center] at (2.8,0.75)
    {$v\mapsto\zeta^{-1}v$\\$u_{1,r}(\zeta)^{-e}$};
  \node[draw,rounded corners,inner sep=6pt] (f1) at (-2.8,-2.05)
    {first fiber: $Q_r$};
  \node[draw,rounded corners,inner sep=6pt] (f2) at (2.8,-2.05)
    {second fiber: $Q_r$};
  \draw[->] (f2.west) -- node[above] {$g\mapsto P_rg$}
    node[below] {$P_rA_r=A_rP_r$} (f1.east);
\end{tikzpicture}
\caption{The two branches of the balanced node and the torsor fibers on
its normalization, shown schematically.  The cover $w^e=t$ multiplies
boundary monodromy by $e$, and the second branch reverses its orientation.  Multiplication
by $P_r$ identifies the resulting action $u_{1,r}^{-e}$ with
$u_{1,r}$ while commuting with the common Frobenius value $A_r$.}
\label{fig:node-gluing}
\end{figure}
